\clearpage\fancyhead[L]{\small POLARIS.}\hypertarget{native-398-part-lower}{}\label{native-398-part-lower}
\begin{center}\hrule\vspace{4pt}\fontsize{10}{12}\selectfont
V. \textsc{Lindenau} u. \textsc{Bohnenberger} Zeitschrift für Astronomie u. verwandte Wiss.\par
Vol. IV. pp. 119..131. July and August. 1817.\par\vspace{4pt}\hrule\vspace{4pt}\end{center}
\begin{flushright}Göttingen, 18 July 1817.\end{flushright}
\par\indent If I am not mistaken, my dearest friend, I have yet to send you the full set of my observations of Polaris with the twelve-inch \textsc{Reichenbach} circle. I am glad that I can \SourceLeaf{399} now add a series of solstice observations, which this year had the benefit of finer weather than ever. These observations were made at the new observatory. As you know, the building is still not quite finished, and the observatory will only really begin operating once the new meridian instruments have arrived and been set up. According to my geodetic determination, the center of the new observatory lies $6''11$ south and $1^{\mathrm s}90$ east of the center of the old one. The spot where I observed Polaris lies $0''12$ north of the center, and the spot for the solar observations lies $0''16$ south of it.

\par\indent You will remember that in some of my earlier letters I mentioned a conjecture about the much-discussed difference between the altitude of the pole obtained from observations of circumpolar stars and that obtained from observations of the Sun. With the smaller repeating circles, the cause might be flexure at the objective end of the telescope, since a heavy brass ring is slipped onto the objective as a counterweight for both the prismatic eyepiece and the clamp. If this were to cause harmful flexure, \textsc{Reichenbach} advised reducing the slip-on weight until it just balanced the telescope properly, and balancing the clamp with a weight screwed onto the alidade circle. For my purposes, however, it was better still to remove the slip-on weight altogether and achieve full balance by screwing a weight onto the alidade circle. If the conjecture had been right, the remaining flexure at the eyepiece end of the telescope would clearly have caused an error in the opposite direction. If flexure had previously made all the zenith distances too small, they would now all have to be too large, and the observations of Polaris would give too low an altitude of the pole and those of the Sun too high an altitude. The result gave no support to this. The altitude of the pole determined from Polaris at the new observatory with the screwed-on weight is almost exactly what one obtains by transferring the altitude observed at the old observatory with the slip-on weight. The solar observations differ in the same direction and by the same amount as before. At first, before the weather allowed a continuous series of astronomical observations, I tried a different way of testing the conjecture. I made several series of observations of the zenith distance of a terrestrial object, some with the slip-on weight and the others with the screwed-on weight. The initial sets of observations did give a difference in the direction and roughly of the size that would have been expected if that really had been the cause. In the later series, however, the difference almost entirely disappeared. In general, the observations of the terrestrial object agreed less closely than the test required, since the object could not be sighted as precisely as objects in the sky.

{\rightskip=0pt plus 1em \par\indent If this exceedingly remarkable difference is due to the instrument, there must be something in its construction that makes us observe zenith distances that are too small. The weight was not the cause. After careful examination, I also found that several other circumstances I was inclined to suspect had no effect. I admit that so far I have found nothing to which I could attribute \emph{such} an effect. An effect in the \emph{opposite} direction would be much easier to explain, and it can indeed occur if one is not on one’s guard. As you know, I experienced this myself last winter. I had taken the circle apart completely, cleaned it, and put it back together. My first observations of Polaris repeatedly gave me an altitude of the pole that was too low by $20''$. Something, then, was increasing the zenith distances. After several other tests I soon found the source of the error in a bend in the eyepiece tube, which is held firmly \SourceLeaf{400} at its far end inside the telescope by the spring. The spring may have lost some of its strength. Once it had been bent a little to restore the necessary tension, the error disappeared at once. The same error will occur if the thread against which the zenith distances are observed is not properly taut. The thread then sags in the middle under its own weight, and the zenith distances observed are evidently too large. It hardly needs saying that the errors at issue are small. If the slack thread sags by a couple of seconds in the middle while its full length may span more than a degree, the eye still detects no curvature unless an image of a perfectly straight horizontal line is projected onto the thread from a suitable distance.\par}

\par\indent It is certain, however, that to declare the instrument free of error we would have to assign the opposite error to other instruments, which show either no such difference between observations of stars and of the Sun or a smaller one. If the 12-inch circle is right, the three-foot circles with upright columns and \textsc{Bessel}’s circle give zenith distances that are too large. If the three-foot circle is right, the twelve-inch circle gives zenith distances that are too small and \textsc{Bessel}’s gives distances that are too large. If the latter is right, all \textsc{Reichenbach} circles give zenith distances that are too small. But no instrument can be judged right with \emph{certainty} unless the possible source of error in the other instruments has been shown definitely and beyond dispute, or the results found with the other instruments have been shown \emph{decisively} to be wrong on other grounds. As things stand, I think we must still admit that no definite judgment can be given. Each astronomer may weigh the relative \emph{probability} differently on subjective grounds. Indeed, it seems natural to me that, when no specific fault can be found with any of the instruments, one feels this most keenly about the instrument one observes with oneself and therefore knows best.

\par\indent There are two further remarks I cannot help adding here.

\par\indent \emph{First}, it seems to me that the testimony of the twelve-inch circles cannot be rejected outright \emph{on the ground} \emph{that} they are the smallest of these instruments, unless it can be shown that they produce constant errors that always act in the same direction \emph{because of their small size}. Their small size might be more likely to lead one to expect random errors that act now one way, now another. Yet experience shows that their results agree just as closely as those of the larger instruments, or at least not much less closely. I do not know whether, \emph{in general}, the size of larger instruments may require us to be at least as much on our guard against constant errors.

\par\indent \emph{Second}, I admit that, leaving all other circumstances aside, an instrument does not yet seem to me to have a presumption in its favor \emph{on the ground} \emph{that} it shows no difference at all between observations of stars and of the Sun. It is still not quite settled whether refraction is \emph{exactly} the same for the Sun and the fixed stars. There is also another important consideration, which I discussed with you several years ago, as you will remember. I raised it then as a possible explanation of the puzzle, without wishing to attach much weight to it. I return to it here because, as far as I know, it has yet to be discussed publicly. We \emph{observe} the center of the Sun’s figure. Its attraction, however, will act on the planets approximately as if it were concentrated at its center of gravity. The Sun’s center of gravity will lie in the plane of the Earth’s orbit, whereas the center of its figure will lie outside that plane. The two cannot be assumed to coincide with mathematical exactness, since the body of the Sun is undoubtedly not a \emph{completely} homogeneous mass. If the center of gravity lay a little south of the center of the figure, the latter would always have \SourceLeaf{401} a small northern latitude. Insofar as this is ignored, observations of the Sun in the northern hemisphere of the Earth would always yield too low an altitude of the pole. Easy and natural though this explanation may seem, it still has its difficulties. It would first require a deeper investigation to see to what extent the body of the Sun could be constituted so that its center of gravity lay south of the center of its figure by $\frac{1}{100}$ or $\frac{1}{1000}$ of the Sun’s diameter without the surface of the Sun, under the conditions of equilibrium, departing more from a sphere than observations show. The observed transits of Mercury must also be considered here. The latitudes of Mercury derived from these transits rest on the assumption that the center of the Sun’s figure lies in the plane of the ecliptic itself, so all the observed latitudes would be too far south. The longitude of the ascending node derived from them would be too large, that of the descending node too small, and the two would fail to agree. The numbers you gave in your theory of Mercury do point to something of this sort, but they do not seem to be compatible with a difference as large as the twelve-inch circles give. Admittedly, too few transits through the descending node have been observed. And since you have given no details of the observations you selected or of how the latitudes were inferred from them, one cannot judge what degree of accuracy to assign to each latitude. It would, however, surely be worth examining the transits of Mercury again with the utmost care from this point of view.

\par\indent It would be extremely interesting if we had observations of the Sun and of the southern circumpolar stars made with repeating circles in the southern hemisphere of the Earth. The southern latitude would, however, have to be high enough that refraction could introduce no appreciable uncertainty for the circumpolar stars.

\par\indent I also take this opportunity to note that Prof. \textsc{Schumacher} sent me the altitude of the pole observed by his brother, Captain \textsc{Schumacher}, at Friedrichsmark, with a \textsc{Reichenbach} astronomical theodolite very like yours. It came to $=55^\circ\,58'\,41''3$ from observations of Polaris and $=55^\circ\,58'\,36''6$ from observations of the Sun. The difference is almost exactly the same as in my observations.

\Needspace{16\baselineskip}\par\indent\emph{Observations of Polaris at lower culmination at the new Göttingen observatory.}

\begin{center}
\fontsize{9}{11}\selectfont
\setlength{\tabcolsep}{2pt}
\begin{tabular}{l|c|c|c|c|c|c}
1817 & \begin{tabular}{c}No. of\\obs.\end{tabular}
& \begin{tabular}{c}Apparent zenith\\distance\end{tabular}
& Barometer & \begin{tabular}{c}Thermo-\\meter\end{tabular}
& \begin{tabular}{c}Re-\\fraction\end{tabular}
& \begin{tabular}{c}Refr. from\\Bessel\end{tabular} \\
\hline
March 28 & 16 & $40^\circ\,7'\,22''40$ & $27^{\mathrm z}\,8^{\mathrm l}2$ & $-0^\circ5$ & $50''78$ & $50''25$ \\
April\ \ 1 & 22 & $24.98$ & $14.9$ & $+3.0$ & $50.88$ & \\
\ \ \ \ \ \ \ \ 2 & 24 & $23.38$ & $13.3$ & $+3.7$ & $50.49$ & \\
\ \ \ \ \ \ \ \ 8 & 28 & $26.36$ & $7\;2$ & $+4.4$ & $49.40$ & $48.87$ \\
\ \ \ \ \ \ \ 11 & 26 & $25.00$ & $9.3$ & $-1.4$ & $51.16$ & \\
\ \ \ \ \ \ \ 13 & 12 & $29.83$ & $8.5$ & $0$ & $50.68$ & \\
May\ \ \ 7 & 30 & $34.68$ & $10.9$ & $+7.5$ & $49.22$ & $48.67$ \\
\end{tabular}
\end{center}

\Needspace{12\baselineskip}\par\indent This gives

\begin{center}
\fontsize{9}{11}\selectfont
\setlength{\tabcolsep}{3pt}
\begin{tabular}{l|c|c|c}
1817 & True zenith dist. & Apparent polar dist. & Altitude of the pole \\
\hline
March 28 & $40^\circ\,8'\,13''18$ & $1^\circ\,40'\,2''51$ & $51^\circ\,31'\,49''33$ \\
April\ \ 1 & $15.86$ & $3.70$ & $47.84$ \\
\ \ \ \ \ \ \ \ 2 & $13.87$ & $4.00$ & $50.13$ \\
\ \ \ \ \ \ \ \ 8 & $15.76$ & $5.79$ & $50.03$ \\
\ \ \ \ \ \ \ 11 & $16.12$ & $6.97$ & $50.51$ \\
\ \ \ \ \ \ \ 13 & $20.51$ & $7.26$ & $46.75$ \\
May\ \ \ 7 & $23.90$ & $13.50$ & $49.67$ \\
\end{tabular}
\end{center}

\SourceLeaf{402}
\par\indent Thus the mean from 158 observations is $51^\circ\,31'\,49''42$, or, reduced to the center of the observatory,

\begin{center}\fontsize{10}{12}\selectfont $51^\circ\,31'\,49''30$ [$-0.644\,Q$ Handwritten note.]\end{center}

\par\indent The refraction is taken from \textsc{Carlini}’s tables, and the polar distance from \textsc{Bessel}’s. My observations of Polaris at lower culmination in 1813 at the old observatory, computed from the same tables, give me its altitude of the pole as $51^\circ\,31'\,55''48$. Thus, reduced to the new observatory, 136 observations give

\[ 51^\circ\,31'\,49''37 \]

\par\indent Changing the counterweight therefore had no effect at all on the observations.

\par\indent I have not yet been able to observe Polaris at upper culmination. On 19 February I observed 6 zenith distances 2 hours away from upper culmination. I will add their calculation here as an example, to fulfill the promise I made in the M.\,C., 1813, p. 484. I am pleased to see from the latest issue of your journal that Mr. \textsc{Littrow}, too, has used this method of observation to good effect.

\Needspace{11\baselineskip}\par\indent My observation times by the clock were

\begin{center}\fontsize{10}{12}\selectfont\setlength{\tabcolsep}{3pt}\begin{tabular}{ccc}
$2^{\mathrm u}$ & $54^{\mathrm m}$ & $19^{\mathrm s}5$\\
 & $59$ & $55.0$\\
$3$ & $4$ & $57.0$\\
 & $8$ & $36.0$\\
 & $13$ & $24.0$\\
 & $15$ & $31.0$\\
\end{tabular}\end{center}

\par\indent (The long intervals were due to the unfavorable weather, which allowed stars to be seen for only a few moments throughout the winter). The arc traversed was $222^\circ\,23'\,53''$, the barometer $27^{\mathrm z}\,10^{\mathrm l}3$, the thermometer $+4^\circ5$, and the star’s time of culmination $0^{\mathrm u}\,55^{\mathrm m}\,23^{\mathrm s}59$. I will explain my method of calculation here only by showing how it is used.

\par\indent The mean of all the clock times is $3^{\mathrm u}\,6^{\mathrm m}\,7^{\mathrm s}08$. The differences from it, disregarding the sign, are

\begin{center}\fontsize{10}{12}\selectfont\setlength{\tabcolsep}{3pt}\begin{tabular}{cc}
$11^{\mathrm m}$ & $47^{\mathrm s}58$\\
$6$ & $12.08$\\
$1$ & $10.08$\\
$2$ & $48.92$\\
$7$ & $16.92$\\
$9$ & $23.92$\\
\end{tabular}\end{center}

\par\indent With these differences I enter the well-known correction table for altitudes near the meridian (I use Mr. von \textsc{Zach}’s table in the Attraction des montagnes) and obtain

\begin{center}\fontsize{10}{12}\selectfont\setlength{\tabcolsep}{3pt}\begin{tabular}{c}
$273''01$\\
$75.50$\\
$2.68$\\
$15.56$\\
$104.11$\\
$173.42$\\
\end{tabular}\end{center}

\noindent Their mean is $107''38$, to which the argument $7^{\mathrm m}\,23^{\mathrm s}73$ corresponds in the same table. I infer from this (by a small trick whose rationale one can easily supply oneself or look up in the Götting. gel. Anz., 1815, p. 454 [23 March, see the reports further below in this volume] \fancyhead[L]{\small SOLAR OBSERVATIONS.}\SourceLeaf{403}) that the mean of the \emph{computed} zenith distances for the \emph{six} actual observation times is approximately equal to the mean of those for the \emph{two} clock times

\begin{center}\fontsize{9}{11}\selectfont\setlength{\tabcolsep}{3pt}\begin{tabular}{ccccccccccc}
$3^{\mathrm u}$ & $6^{\mathrm m}$ & $7^{\mathrm s}08$ & $-$ & $7^{\mathrm m}$ & $23^{\mathrm s}73$ & $=$ & $2^{\mathrm u}$ & $58^{\mathrm m}$ & $43^{\mathrm s}35$\\
$3$ & $6$ & $7{,}08$ & $+$ & $7$ & $23{,}73$ & $=$ & $3$ & $13$ & $30{,}81$\\
\end{tabular}\end{center}

\Needspace{12\baselineskip}\noindent to which the hour angles $30^\circ\,49'\,56''9$ and $34^\circ\,31'\,48''3$ correspond. The zenith distances computed by the familiar formulas for these hour angles, using the declination $88^\circ\,20'\,7''53$ taken from \textsc{Bessel}’s table and the altitude of the pole $51^\circ\,31'\,47''0$, are

\begin{center}\fontsize{10}{12}\selectfont\begin{tabular}{lrrr}
& $37^\circ$ & $2'$ & $57''31$\\
& $37$ & $6$ & $32.75$\\\cline{2-4}
giving the mean & $37$ & $4$ & $45.03$\\
\end{tabular}\end{center}
\noindent or, allowing for the corrections $d\delta$, $d\varphi$ to the adopted declination and altitude of the pole\footnote{The coefficients were taken using the mean hour angle from a table computed by Mr. \textsc{Westphal} for the local altitude of the pole.}

\[ = 37^\circ\,4'\,45''03 - 1.00\,d\varphi + 0.83\,d\delta \]

\par\indent The mean of the observed zenith distances is $37^\circ\,3'\,58''83$, or, with the refraction $44''64$ removed,

\[ = 37^\circ\,4'\,43''47 \]

\par\indent It follows that $d\varphi = +1''56 + 0.83\,d\delta$, or the altitude of the pole $= 51^\circ\,31'\,48''56 + 0.83\,d\delta$. The observations at lower culmination yielded $51^\circ\,31'\,49''42 - d\delta$, from which $d\delta = +0''47$ would follow, and an altitude of the pole of $51^\circ\,31'\,48''95$. (Of course I have added this only by way of example. To make the altitude of the pole independent of the adopted declination, the observations near upper culmination would have to be much more numerous, and the results would then have to be combined by the method of least squares). One could also introduce a correction to the adopted right ascension, but it is better to leave it out, since this element is now known so precisely thanks to your efforts and those of our \textsc{Bessel}.

\Needspace{22\baselineskip}\par\indent I now give my observations of the Sun.

\begin{center}\fontsize{9}{11}\selectfont\setlength{\tabcolsep}{3pt}\begin{tabular}{c|c|c|c|c|c}
$1817$ & \shortstack{Arc\\traversed} & \shortstack{Sum of the\\reduction} & \shortstack{No. of\\obs.} & Barometer & Thermometer\\
\hline$\text{June}\ 13$ & $283^\circ\,11'\,9''75$ & $11'\,28''90$ & $10$ & $27^{\mathrm z}\,6^{\mathrm l}6$ & $20^\circ8$\\
$16$ & $451\;21\;38.75$ & $52\,1.70$ & $16$ & $13.1$ & $13.0$\\
$17$ & $563\;46\;30.75$ & $77\,32.11$ & $20$ & $11.0$ & $14.3$\\
$18$ & $449\;59\;10.75$ & $25\,45.01$ & $16$ & $7.7$ & $16.4$\\
$19$ & $337\;8\;25.25$ & $12\,31.98$ & $12$ & $7.7$ & $19.3$\\
$20$ & $562\;18\;1.50$ & $61\;30\;15$ & $20$ & $7.4$ & $17.0$\\
$21$ & $337\;0\;25.50$ & $19\,20.42$ & $12$ & $7.5$ & $21.0$\\
$22$ & $280\;43\;13.25$ & $7\,55.97$ & $10$ & $8.3$ & $22.0$\\
$23$ & $337\;14\;3.75$ & $26\,26.97$ & $12$ & $8.1$ & $20.7$\\
$24$ & $449\;48\;49.25$ & $31\,1.55$ & $16$ & $8.1$ & $22.5$\\
$25$ & $281\;19\;53.25$ & $18\,40.57$ & $10$ & $8.1$ & $20.4$\\
$26$ & $394\;13\;39.25$ & $23\,47.90$ & $14$ & $7.5$ & $19.0$\\
$27$ & $620\;53\;26.75$ & $73\,42.54$ & $22$ & $5.8$ & $17.5$\\
$29$ & $452\;46\;31.50$ & $40\,2.88$ & $16$ & $8.9$ & $16.0$\\
$30$ & $226\;51\;23.75$ & $21\,7.83$ & $8$ & $7.6$ & $19.0$\\
\end{tabular}\end{center}

\SourceLeaf{404}
\Needspace{23\baselineskip}\par\indent From this I further find

\begin{center}\fontsize{9}{11}\selectfont\setlength{\tabcolsep}{3pt}\begin{tabular}{c|c|c|c|c|c}
$1817$ & Refraction & Parallax & Latitude & \shortstack{Reduction to\\the solstice} & \shortstack{Zenith distance\\at the solstice}\\
\hline$\text{June}\ 13$ & $29''22$ & $4.01$ & $-0''11$ & $14'\,30''92$ & $28^\circ\,3'\,52''26$\\
$16$ & $30.73$ & $3.99$ & $+0.26$ & $5\,59.44$ & $48.63$\\
$17$ & $30.33$ & $3.99$ & $+0.34$ & $3\,58.31$ & $55.30$\\
$18$ & $29.67$ & $3.99$ & $+0.38$ & $2\,21.96$ & $53.46$\\
$19$ & $29.25$ & $3.99$ & $+0.39$ & $1\,10.43$ & $54.66$\\
$20$ & $29.17$ & $3.99$ & $+0.37$ & $0\,23.71$ & $51.41$\\
$21$ & $28.99$ & $3.99$ & $+0.30$ & $0\,1.81$ & $48.91$\\
$22$ & $28.91$ & $3.99$ & $+0.20$ & $0\,4.75$ & $52.10$\\
$23$ & $29.11$ & $3.98$ & $+0.08$ & $0\,32.50$ & $51.77$\\
$24$ & $28.89$ & $3.98$ & $-0.06$ & $1\,24.96$ & $51.62$\\
$25$ & $29.12$ & $3.99$ & $-0.22$ & $2\,42.22$ & $49.96$\\
$26$ & $29.35$ & $3.99$ & $-0.37$ & $4\,24.19$ & $51.61$\\
$27$ & $29.47$ & $3.99$ & $-0.52$ & $6\,30.75$ & $53.49$\\
$29$ & $30.05$ & $4.01$ & $-0.79$ & $11\,57.68$ & $51.86$\\
$30$ & $29.60$ & $4.01$ & $-0.88$ & $15\,17.82$ & $53.88$\\
\end{tabular}\end{center}

\par\indent Mean of 214 observations

\[ = 28^\circ\,3'\,52''15 \]

\par\indent My theory of probability gives me the probable error of this result $= \pm 0''32$. According to the Milan ephemerides (from which Captain \textsc{Schumacher} also took the solar declinations for reducing his observations mentioned above), the apparent obliquity of the ecliptic is $= 23^\circ\,27'\,52''9$, so the altitude of the pole is $51^\circ\,31'\,45''05$, or, reduced to the center,

\begin{center}\fontsize{10}{12}\selectfont $50^\circ\,31'\,45''21$ [$+0.472\,Q$ Handwritten note.]\end{center}

\noindent which is $4''09$ less than the value from Polaris. With \textsc{Bessel}’s obliquity of the ecliptic, the difference would be another second and a half larger.

\par\indent I should also mention that the reading of the thermometer at the barometer was not recorded. At such moderate zenith distances, neglecting the reduction to a standard temperature or to the temperature of the air outdoors can make a difference of only a very small fraction of a second. If anyone wishes to compute the observations using other refraction tables and allowing for such a refraction, they will be little in error if, in each case, they assume the mercury to be $2^\circ$ colder than the air outdoors for the solar observations and $2^\circ$ warmer for the observations of Polaris. I am prompted to make this remark by the very well-founded reminder of Prof. \textsc{Pasquich} in the new issue of the journal that has just been handed to me. [Handwritten: $51^\circ\,30'\,48''26 - 0.966\,Q$ Winter solstice 1817.]
